\documentclass[10pt]{article}
\usepackage[margin=20mm]{geometry}
\usepackage{amsmath,amssymb,mathtools,fontspec,unicode-math,xurl,hyperref}
\IfFontExistsTF{DejaVu Serif}{\setmainfont{DejaVu Serif}}{\setmainfont{Latin Modern Roman}}
\setmathfont{latinmodern-math.otf}
\setlength{\parindent}{0pt}\setlength{\parskip}{6pt}\sloppy
\title{A verification of selected arguments in
Affirmative Resolution of Bourgain’s Slicing Problem
Using Guan’s Bound}
\author{}
\date{}
\begin{document}
\maketitle
Abstract: The final implication from the stated covariance estimate to a universal isotropic-constant bound and the slicing assertion is checked, conditional on the paper’s geometric and entropy-stability inputs. The matrix order, dimension dependence and endpoint in the entropy argument are consistent. Equation (56) contains a local extraneous square, corrected by the following line of the source. One Minor finding is supported; no substantive defect in the checked implication is established.\par
\section*{Identity and normalizations}
The source is Boaz Klartag and Joseph Lehec, Affirmative resolution of Bourgain’s slicing problem using Guan’s bound, arXiv:2412.15044v1, 19 December 2024, 20 pages [KL]. The assigned DOI is 10.1007/s00039-025-00718-w. Equivalence with the journal version has not been established. Page references use the acquired PDF’s physical and printed numbering.\par
An isotropic probability measure is centered with covariance \(\mathrm{Id}\). For a volume-one convex body \(K\subset\mathbb R^n\), \(L_K=(\det\operatorname{Cov}(K))^{1/(2n)}\), and \(L_n\) is the supremum over such bodies. The entropy convention (3), p. 3, gives \(L_\mu=\exp(-\operatorname{Ent}(\mu)/m)\) for an isotropic measure on \(\mathbb R^m\). Theorem 1.1 requires some affine hyperplane section, with ordinary Euclidean \((n-1)\)-dimensional volume.\par
\section*{Problem 1: the derivative in equation (56)}
Assessment: Minor, local formula error. On p. 15 the final term of (56) is printed as \[(1-r)^{-3}\frac{d}{dt}\mathbb EA_t^2 .\] Since \(\Gamma_r=A_t/(1-r)\) and \(t=r/(1-r)\), the chain rule requires \[\frac{d}{dr}\mathbb E\Gamma_r
 =\frac{\mathbb E\Gamma_r}{1-r}
  +\frac{1}{(1-r)^3}\frac{d}{dt}\mathbb EA_t .\] This distinction is substantive for the displayed equality: for standard Gaussian initial data, \(A_t=(1+t)^{-1}\mathrm{Id}\) and \(\Gamma_r=\mathrm{Id}\). The derivative is zero, whereas the printed right-hand side at \(r=1/4\) equals \(-2\mathrm{Id}/3\).\par
Deleting the superscript \(2\) inside the derivative repairs (56). The next displayed line already uses \(d(\mathbb EA_t)/dt=-\mathbb EA_t^2\) and obtains the correct Lemma 3.3(iv). Thus this finding does not refute that lemma or the main theorem. The assessment concerns only the acquired version.\par
\section*{Conditional verification of the entropy argument}
\subsection*{The lower bound for the mean covariance}
Let \(C_G\) be the universal constant in Lemma 2.3. The identities \(d(\mathbb E\operatorname{Tr}A_t)/dt=-\mathbb E\operatorname{Tr}A_t^2\) and \(\mathbb EA_t\le\mathrm{Id}\), with \(t_0=\min\{1/(2C_G),1\}\), give \(\mathbb E\operatorname{Tr}A_{t_0}\ge n/2\). The eigenvalue selection on p. 10 and conditional-covariance monotonicity in Lemma 2.2 give the covariance part of Proposition 2.4. Its isotropic-constant transfer uses the separate geometric Lemma 2.5.\par
Accordingly, for \(n\ge10\), take the proposition’s isotropic measure \(\mu\) on \(\mathbb R^m\), where \(n/4\le m\le n\), \(L_\mu\ge c_pL_n\) and \(\mathbb EA_{t_0}\ge\mathrm{Id}/4\). Here \(c_p>0\) is universal. Put \[\xi=\frac{t_0}{1+t_0},\qquad M(r)=\mathbb E\Gamma_r,\qquad
 a=\frac{\xi}{4},\qquad q(r)=\mathbb E|v_r|^2.\] Lemma 3.3 gives \(\Gamma_r^2\le\Gamma_r/r\), hence \(M'(r)\ge-M(r)/r\). Thus \(rM(r)\) is increasing and, for \(r\ge\xi\), \[a\mathrm{Id}\le \frac{\xi}{r}M(\xi)\le M(r)\le\mathrm{Id}.\] Both \(\xi\) and \(a\) are positive universal constants.\par
\subsection*{The variance budget and matrix order}
Write \(\|\cdot\|_{\mathrm{HS}}\) for the Hilbert–Schmidt norm. Since \(Q(r):=\mathbb E(v_r\otimes v_r)=(\mathrm{Id}-M(r))/(1-r)\ge0\), positive-semidefinite trace order gives \[\frac{\|\mathrm{Id}-M(r)\|_{\mathrm{HS}}^2}{1-r}
 =\operatorname{Tr}\bigl(Q(r)(\mathrm{Id}-M(r))\bigr)\le(1-a)q(r).\] No commutativity of random matrices is used. In Lemma 3.4, the bound \(\Gamma_r\le\xi^{-1}\mathrm{Id}\) makes the matrix denominator at most \(2\xi^{-1}\mathrm{Id}\); independence of its two copies gives the stated variance coefficient. Together with Lemma 3.5, this yields \[V:=\int_\xi^1
 \frac{\mathbb E\|\Gamma_r-M(r)\|_{\mathrm{HS}}^2}{1-r}\,dr
 \le \frac{8m}{\xi}.\]\par
\subsection*{The endpoint in (64)}
The endpoint can be justified before assuming the desired entropy bound. For \(R<1\), let \(I_R=\int_\xi^R q(r)\,dr\). Lemma 3.3(iii) gives \(q'\ge0\), and ordinary integration by parts gives \[\int_\xi^R(1-r)q'(r)\,dr
 =I_R-(1-\xi)q(\xi)+(1-R)q(R).\] Using the variance decomposition (66) and the trace estimate above, \[aI_R\le V+(1-\xi)q(\xi).\] The omitted endpoint term is nonnegative. This bound first proves \(\int_\xi^1q<\infty\); then \(0\le(1-R)q(R)\le\int_R^1q\to0\). It therefore justifies (64) without a circular assumption.\par
\section*{Constants and the slicing conclusion}
Lemma 3.1 gives \(q(\xi)\le4m/(1-\xi)^2\). Consequently, \[\int_\xi^1 q(r)\,dr
 \le \left(\frac{32}{\xi^2}
          +\frac{16}{\xi(1-\xi)}\right)m .\] Monotonicity of \(q\) gives \(\int_0^1q\le(1-\xi)^{-1}\int_\xi^1q\). By Corollary 3.2, \[\frac{D(\mu\Vert\gamma_1)}m\le B,\qquad
 B=\frac{16}{\xi^2(1-\xi)}
   +\frac{8}{\xi(1-\xi)^2}.\] The entropy normalization therefore yields \[L_n\le c_p^{-1}L_\mu
       \le \frac{c_p^{-1}e^B}{\sqrt{2\pi e}}
       \qquad(n\ge10).\] This is dimension-free; all entropy expressions use the selected dimension \(m\). The finitely many smaller dimensions are absorbed in the constant. No numerical value is asserted for the universal input constants.\par
For completeness, the required Euclidean section conclusion follows directly. Let \(X\) be uniform on a volume-one body \(K\), and choose a unit eigenvector \(u\) for the least eigenvalue \(\sigma^2\) of its covariance. Then \[\sigma^2\le(\det\operatorname{Cov}(K))^{1/n}=L_K^2.\] By coarea, the density of \(u\cdot X\) at \(s\) equals the section volume of \(K\cap\{u\cdot x=s\}\). Chebyshev’s inequality places at least \(3/4\) of this marginal mass in an interval of length \(4\sigma\). Hence some section has volume at least \(3/(16\sigma)\ge3/(16L_K)\). A uniform bound \(L_n\le C_*\) gives Theorem 1.1 with, for example, \(c=3/(32C_*)\), allowing its strict inequality. The one-dimensional case is immediate. No affine invariance of Euclidean section area is assumed.\par
\section*{Evidence, limits and conclusion}
The full acquired text and manifest were read. The detailed checks above concern the final implication, its constants, matrix inequalities, endpoint and normalization; full reading is not a proof audit of every preliminary result. Guan’s original covariance proof, the geometric inputs behind Lemma 2.5, and the original Eldan–Mikulincer and Ball–Nguyen entropy results were not independently checked from their source papers. They remain named inputs to the conditional verification.\par
The primary executor separately checked the acquired text and the constant, endpoint and Gaussian calculations. These calculations are analytic; no author code or numerical experiment is used as a substitute for proof. Compilation and report-page inspection are recorded separately.\par
Problem 1 is local and does not impair the subsequent calculation. No further substantive defect is established within this checked scope. Scientific decision remains pending, human review entries are empty, and independent acceptance is incomplete.\par
1 B. Klartag and J. Lehec, Affirmative resolution of Bourgain’s slicing problem using Guan’s bound, arXiv:2412.15044v1, 19 December 2024.\par
\end{document}
